\chapter{Without Grammar}

Imagine a world where language is stripped down to its bare essentials, where communication is a dance of simple sounds and gestures, unencumbered by the complexities of grammar. This is the world of a child on the cusp of language discovery, a world where every word is an exploration and every sound a new frontier.

Ethan was 11 months and 21 days old when some researchers visited him to record him and his mother. At this tender age, Ethan was just getting started at speaking.

%https://sla.talkbank.org/TBB/phon/Eng-NA/Providence/Ethan/001121.cha

\begin{description}
    \item[Mother:] Should we find your ball? Where's your ball? Yes, whoops sorry where's your ball? Let's see. Uhoh!
    \item[Ethan:] Ball.
    \item[Mother:] Yeah it's your ball. Bounce, bounce, bounce. Ready? Bounce, bounce, bounce, bounce whoops!
    \item[Mother:] It hit you in the head. Bounce, bounce, bounce. Bounce, bounce, bounce, bounce, bounce, bounce, bounce, bounce, bounce, bounce, bounce! You wanna see the pumpkin again?
    \item[Ethan:] Baby.
    \item[Mother:] Baby. That's all the pictures. Baby.
    \item[Ethan:] (laughs)
    \item[Mother:] That's a pumpkin. Want me to turn the page? Who's that?
    \item[Ethan:] That.
    \item[Mother:] Baby. \&-uh, there. You're sleepy. You're a sleepy boy.
\end{description}

I think it's fair to say that Ethan's language is ungrammatical, just in a slightly different sense than what how we usually use the word. When we say something is ungrammatical, typically we mean that it fails to conform to existing grammatical norms, but it could also mean that no such norms exist. Perhaps Ethan's case is more like the second kind of ungrammatical. His speech comes out in single words: \data{ball}, \data{baby}, \data{that}. In other parts of the interaction, he says, \data{ha cow} or \data{quack quack}, and \data{ha bus}, but there's no grammatical connection between any of the words.

Ethan's world is one we've all been through. Ernst Haeckel once said, ``ontogeny recapitulates phylogeny.'' He meant the development of an embryo echoes the development of the species, going all the way back to a single-celled ancestor. Though it turned out that Haeckel was a bit wide of the mark, there was something to the idea. And we can see that same something at work in language learning. A child has to perceive meaning in the language around it if it's to figure out the grammar. %is this true?

This example from early \textit{English child language}, discussed by \autocite{Jackendoff2014}.

\ea[]{English child language (Brown 1973, quoting Lois Bloom)\\ \textit{Mommy sock}}
    \ea[]{`Mommy is putting a sock on me’}
        \ex[]{`Mommy's sock’}
    \z
\z

The pair of words \data{Mommy} + \data{sock} is used at different times by the child to mean either (a) `Mommy is putting a sock on me' or (b) `Mommy's sock'. There is nothing in the grammar that says which meaning is which or how this should be composed. 

\begin{quote}
    it would seem reasonable to posit a stage in early child language where words are strung together without benefit of any kind of notion of predication \autocite[62]{Gil2014}
\end{quote}

(what about great apes. Gil also mentions those. See Hudson 2012 too)

In fact, the ability to discern meaning before understanding grammar is a fundamental aspect of language learning and evolution. Consider an English speaker in Italy hearing the sentence \data{Il museo d'arte moderna è interessante}. Even without knowledge of Italian grammar, the speaker might recognize familiar words and work out the meaning: `The modern art museum is interesting.'

This process of extracting meaning from unfamiliar language structures reflects a broader pattern in both individual language acquisition (linguistic ontogeny) and the historical development of individual languages (linguistic phylogeny), along with the birth of human language itself (what I'll call \textsc{logogeny}). 

\section{Learning a language}

    In the early stages of learning a new language, a learner often relies on recognizable words, context, and non-verbal cues to grasp the meaning of a sentence. The grammar, though present, might remain opaque until a later stage of learning.

\section{The birth of human language (logogeny)}

Similarly, in the evolution of human language, the development of meaning likely preceded the development of complex grammatical rules. Early forms of communication might have consisted of simple lexical items and basic semantic relationships, much like the initial stages of language learning. Words and symbols were created to represent objects, actions, and concepts, forming the foundational building blocks of communication.

As the lexicon grew and the need for more nuanced expression arose, grammatical rules and structures began to evolve. Syntax and morphology became more complex, allowing for the expression of more intricate ideas and relationships. The transition from simple lexical items to a fully developed language system was gradual, guided by the interplay between meaning and grammar.

In \textit{The development of meaning} and \textit{The development of grammar}, James Hurford posits a two-stage development process. In the first stage, the development of meaning takes precedence. This stage is characterized by the formation of simple lexical items and the establishment of basic semantic relationships. Words and symbols are created to represent objects, actions, and concepts, forming the foundational building blocks of communication.

In the second stage, the development of grammar emerges. Once a sufficient lexicon has been established, grammatical rules and structures begin to evolve. Syntax and morphology become more complex, allowing for the expression of more nuanced ideas and relationships. This stage marks the transition from a rudimentary form of communication to a fully developed language system, complete with its own internal logic and consistency.

Hurford's thesis provides valuable insights into the intricate interplay between meaning and grammar in the evolution of human language. It emphasizes the sequential nature of this development, highlighting the importance of semantic foundations in paving the way for grammatical complexity. This perspective contributes to our understanding of how languages have evolved over time, reflecting the adaptability and creativity of the human mind.


\section{The Continuity of Ungrammatical Communication}

Ethan's world, where language is unbound by grammar, is not confined to infancy. It persists throughout our lives, manifesting in various forms of communication where grammar is simply irrelevant.

\subsection{Pragmatic Communication}

We often engage in communication that relies on context, shared understanding, and non-verbal cues rather than strict grammatical structure. Examples include:

\begin{itemize}
\item[Gestures]: Simple nods or waves that convey meaning without words.
\item[Onomatopoeia]: Sounds that imitate the natural noises they represent.
\item[Emoji]: Symbols that convey emotions and ideas without traditional syntax.
\end{itemize}

\subsection{Creative Language Use}

Ungrammatical communication can also be a creative tool:

\begin{itemize}
\item[Poetry and Art]: Breaking grammatical rules to evoke emotions or create aesthetic effects.
\item[Advertising]: Ignoring grammatical norms to create memorable phrases.
\end{itemize}


\subsection{Conclusion}

The ungrammatical mode of communication that begins in childhood continues to be a vital part of human interaction. It's not a phase to be outgrown but a fundamental aspect of how we connect with each other. Whether through simplicity, creativity, or technical precision, ungrammatical communication transcends the boundaries of traditional language structure, reflecting the multifaceted nature of human expression.

This world is not left behind as we grow; it grows with us, enriching our ability to convey meaning in diverse and profound ways. It's a reminder that language is not merely a set of rules but a living, evolving entity that adapts to our needs and desires.

\section{Riau Indonesian: A Linguistic Enigma}

Sumatra, one of Indonesia's largest islands, is a land of contrasts and diversity. The western part of the island is a region marked by a series of volcanoes and ex-volcanoes, a testament to Sumatra's tectonic activity. These towering peaks, some still active, are interspersed with remnants of craters that have long since become dormant. The volcanic soil here is highly fertile, giving rise to lush vegetation and supporting a diverse array of wildlife.

In stark contrast to the rugged western terrain, the eastern part of Sumatra is characterized by lowlands, swamps, and dense tropical forests. This region is crisscrossed by numerous rivers that nourish the land, supporting a unique ecosystem teeming with various species of birds, reptiles, and amphibians. It is within this rich and diverse landscape that the province of Riau is found, occupying the central eastern part of the island with Jambi province to the south, the province of North Sumatra (Sumatera Utara) to the north, and to the east, across the Malacca Strait, Malaysia and Singapore.

The cultural heritage of Riau is a rich blend of Malay, Chinese, and indigenous influences shaping its traditions, arts, cuisine, and language. Riau's strategic location near the Malacca Strait has historically made it a hub for trade and commerce. This narrow but strategically significant stretch of water serves as the main shipping channel between the Pacific and Indian Oceans. As one of the world's most important and busiest maritime routes, the Malacca Strait plays a vital role in global trade and economics, influencing not only the region but the entire world.

The cultural impact of the Malacca Strait on the province of Riau has been profound. The constant movement of goods, people, and ideas through the strait has led to a rich tapestry of cultural influences in Riau. The Malay people, who have historically inhabited the region, were influenced by traders and settlers from India, China, the Arabian Peninsula, and various European nations.

The influence of Islam, brought by Arab traders, is particularly notable in Riau, shaping religious practices and cultural traditions. Chinese and Indian influences are evident in the local cuisine, arts, and architecture. The colonial legacy of the European powers has also left its mark on the administrative and legal systems of the region.

Riau's ports, such as Dumai and Tanjung Pinang, became melting pots of diverse cultures, fostering a unique blend of traditions and customs. The intermingling of different ethnic groups led to the creation of new languages, art forms, and social practices, contributing to the rich cultural heritage of the province.


The province of Riau , where the Siak River winds its way through dense tropical forests, lies a linguistic phenomenon that challenges our conventional understanding of grammar. Here, the people speak Riau Indonesian, a language that relies a good deal on our abilities to extract meaning from context, and very little on grammar.

Riau Indonesian is not a broken or simplified form of the standard Indonesian language. Nor is it a pidgin or a creole. A pidgin is a simplified language that develops as a means of communication between groups that do not share a common language. Pidgins typically emerge in trade, colonization, or other situations where speakers of different languages need to communicate for practical purposes. 

Because of how they develop, pidgins lack native speakers, or native signers, in the case of signed languages. And just as I, a non-native speaker, struggle to be consistent with the grammar of Italian, a language I've been studying for about eight months, pidgin speakers are often inconsistent in the way they use the language. In other words, pidgins don't have have any particular word order. Not only that, but words rarely have more than one form, no past tense or plural number. And pidgins always come with a at least two different accents built in. In other words, they lack syntax, most morphology, and a regular phonology. They lack grammar.

But Riau Indonesian is clearly not a pidgin. It has been around for a very long time and has many native speakers. Given this, some linguists claimed that it must, then, be a creole. A creole is a relatively young language. Creoles often emerge in situations where a pidgin becomes the first language of a community, such as in enslaved or indentured servant populations where speakers of various languages were brought together. Unlike pidgins, then, creoles have native speakers, folks who have grown up speaking the language. And when a group of people grow up speaking a language together, they tend to speak it in a consistent way. In other words, Creole's typically have more grammar than pidgins do, though they generally have little or no morphology: no difference between singular and plural noun forms, no past-tense verb forms, no comparative adjective forms. Of course they can express plurality and past time and comparison; they just don't do it by changing the form of the word in question.

Some underspecification here and there is normal in any language. In Japanese, for instance, which has never been accused of having a particularly simple grammar, \data{バナナ食べた} (banana tabeta) can mean `I ate a banana,' `I ate the banana,' `did you eat the bananas?' or various other combinations. Japanese doesn't typically mark a distinction between singular `banana' and plural `bananas', or between definite `the banana' and indefinite `a banana'. If this strikes you as strange or confusing, consider that an English speaker wouldn't typically differentiate between definite \data{I ate the} (\data{individual whole}) \data{banana} and \data{I at the banana} (\data{mush}), even though we often mark a unit/mass distinction in indefinite cases, such as \data{I ate a banana} vs \data{I ate some banana} (\data{mush}).

One story we English speakers might tell is that the speaker anticipates that the other person has enough context to know which banana they're talking about when they say \data{the banana}. That's why they use the definite article \data{the}. And so, if you have that context, then you don't need to further point out that it's a unit or a mass. But the Japanese speaker would say that the context is usually sufficient to understand the meaning without needing to mark distinctions between singular and plural or definite and indefinite. And how important is it anyways to distinguish between singular and plural? Old English used to have dual pronouns -- \data{wit}, for instance, was `we two' -- but Modern English speakers don't seem to suffer from the loss of that distinction. No language distinguishes grammatically between seven bananas and eight.

In Riau Indonesian, \data{makan pisang} (`eat' + `banana') can be combined with 

John McWhorter, who cut his teeth in linguistics working on creoles. 

But David Gil s


\textbf{Creoles:}
A creole is a stable, fully developed natural language that has developed from a pidgin. Creoles often arise in situations where a pidgin becomes the first language of a community, such as in enslaved or indentured servant populations where speakers of various languages were brought together. Some features of creoles include:
\begin{itemize}
    \item More complex and stable than a pidgin, with its own grammar and vocabulary.
    \item Often has native speakers, especially in subsequent generations.
    \item Used in a wide range of social contexts, not just specific situations.
    \item May retain elements from the parent languages but has its own unique characteristics.
\end{itemize}




Consider the following conversation in Riau Indonesian:

\begin{description}
    \item[Person A:] You go where?
    \item[Person B:] Market.
    \item[Person A:] Buy what?
    \item[Person B:] Fish.
\end{description}

The sentences are stripped down to their bare essentials, yet they convey complete thoughts. There's no need for conjunctions, articles, or tenses. The words flow naturally, unencumbered by the rules that govern most languages.

This linguistic simplicity is not a sign of lack or limitation. On the contrary, it's a testament to the power of unadorned communication, where meaning is derived not just from the words themselves but from the context, the shared understanding between the speakers, the tone, and the gesture.

\section{Language with less grammar}

Ethan and his mother were speaking in English, and Ethan simply wasn't using the grammatical apparatus of that language. But many languages simply have less grammatical apparatus to be used. One such language, documented 

Riau Indonesian is not an anomaly; it's a mirror reflecting a universal truth about human communication. It reminds us that language is not a rigid construct but a living, breathing entity that evolves to suit our needs. It echoes the world of Ethan, where words are free from grammatical constraints, where communication is a dance of sound and meaning.

In the heart of Sumatra, Riau Indonesian stands as a poetic affirmation of our innate ability to connect with one another, to transcend the boundaries of grammar, and to find understanding in the simplicity of human expression. It's a language that speaks not just to the mind but to the soul, resonating with the very essence of what it means to communicate.

\begin{quotation}
    no reason to believe that speakers of a language such as Riau Indonesian always worry about whether a clause unmarked for [tense] has past, present or future reference any more than speakers of a language such as English invariably puzzle over past-tense-marked clauses, wondering whether they refer to a few hours back, a day ago, a week to a month ago, a couple of months to a couple of years ago, or the distant or legendary past, just because these distinctions happen to be grammaticalized in a language such as Yagua.\footnote{Yagua is an Amazonian language with seven tenses \autocite[385]{Payne1990}.} \hfill\autocite[43--44]{Gil2014}
\end{quotation}

\begin{quotation}
    How to translate (10) into English? The grammar of English presents the translator with a dilemma, a forced choice between one of the two translations given in (10i) and (10ii). Whereas the former takes \textit{cantik} ‘beautiful’ to be predicated of \textit{gol} ‘goal’, in the latter \textit{cantik} is assumed to be attributed of \textit{gol}. These two translations thus point towards distinct potential grammatical analyses, involving distinct semantic and perhaps also syntactic structures. So which of the two translations is the right one—or is the sentence perhaps ambiguous between the two? The right answer is none of the above. Riau Indonesian provides little or no evidence for the grammaticalization of a predicate/ attribute distinction. Moreover, in the actual utterance context, there is no reason to assume that the speaker had in mind one of these two translations to the exclusion of the other. Instead, the sentence is vague or unspecified with respect to the predicate/attribute distinction. It has a single unitary meaning which may be represented in terms of the association operator as A ( GOAL, BEAUTIFUL ), or, in plainer English, ‘something to do with goal and beautiful’. In other words, the semantics of(10) is like that of the single representation in (8), not the two in (9). Examples such as (10), which are plentiful in any naturalistic corpus of Riau Indonesian, show how a language can enjoy an expressive power comparable to those of most familiar languages without recourse to predication. \hfill\autocite[57]{Gil2014}
\end{quotation}

It's not the case that Riau Indonesian is entirely lacking in grammar. A preposition + noun phrase combination such as \data{in} + \data{the world} or \data{from} + \data{Canada} invariably has a preposition like \data{dari} (`from') appearing before the NP. So, \data{dari Kanada} would be grammatical but *\data{Kanada dari} would not \autocite{Gil2000}. %also personal communication Aug 23, 2023. 
There seems to be a general head-precedes-modifier principle \autocite{Gil2005}. 

To explain that idea, I'll start with an example from English. Consider the sentence \data{birds fly}. We can think of this as a noun phrase (NP) \data{birds} followed by a verb phrase \data{fly}. The NP can have a modifier such as the adjective \data{typical}, so that we now have \data{typical birds fly}. Here \data{typical} belongs to the category \textsc{adjective} and performs the function of \textsc{modifier}. But that leaves us with the question, what is the function of \data{birds} in the NP. Linguists call it the \textsc{head} of the phrase.

In English, it's generally true that modifiers precede heads: \data{typical birds fly} follows this pattern, and if we switched it around, it would be the ungrammatical *\data{birds typical fly}. But that's not always the case. English also has NPs like \data{governors general}, where \data{governors} is the head and it's followed by the modifier \data{general}. And in verb phrases (VPs), it's actually quite common for a modifier, such as the adverb \data{gracefully} to follow the head verb, as in \data{birds fly gracefully}.

Now, returning to Riau Indonesian, 

\begin{quotation}
    condensing a mess variety, and sections of this chapter so las, creoles pagins, the B normal innate capac sign languages give us a windedles. or humans with on the ambient hasity for language so about construcing her anguages whelfed, an ust planguages are mixed and inconsistent, of parfacently exemplited, or just plain absent. The most basic characteristics ahistorically new languages are a relative lack of inflectional morphology and hiation words. These traits bring with them some concomitant features, par-siularly the use of single-valence verbs and serial verb constructions, motivated by the lack of morphology and grammatical marking devices. Subordination is hater than in older languages. These new languages find ways, early in their development, of highlighting some sentence elements for discourse purposes, spically by exploiting the front position. The more recent research surveyed here reinforces the vision articulated over thirty years ago by Talmy Givón (19792, p. 296), when he listed the 'major salient properties of the presyntactic, pragmatic mode of discourse' as (1) Topic-comment word order, (2) Concate-natn, (3) Low noun per verb ratio, (4) Lack of grammatical morphology, (5) nionation, and (6) Zero anaphora. 59 Givón was ahead of his time, and these it now ideas whose time has gradually come, to some sooner than to others. \hfill\autocite[459]{Hurford2012a}
\end{quotation}

